% ============================================================
% Case Study FRFP Case Study (Health Systems / Industry):
% IBM Watson Health — Rise and Fall of a Healthcare AI Platform
% ============================================================

\section{FRFP Case Study (Health Systems / Industry): IBM Watson Health}
\label{sec:frfp-watson-health}

This chapter looks beyond Watson for Oncology and analyses the broader arc of
\emph{IBM Watson Health} as an FRFP case study at platform scale. We treat Watson
Health as a large organisational attempt to embed AI products across healthcare,
and use its trajectory to illustrate FRFP concerns around explicit/tacit
separation, boundary design, safe horizons, and institutional governance.

\subsection{Background: From Jeopardy! to Healthcare Platform}

After Watson’s 2011 Jeopardy! success, IBM announced ambitious plans to apply
Watson to healthcare, pitching it as a system that could ingest large volumes of
medical literature and help clinicians make better decisions. By 2015, IBM had
created the Watson Health business unit, acquired multiple health data companies
(e.g.\ Truven Health Analytics, Phytel, Explorys), and pursued high-profile
partnerships with institutions such as MD Anderson Cancer Center, Memorial Sloan
Kettering, Mayo Clinic, and others.

Watson Health’s portfolio included:

\begin{itemize}
  \item Watson for Oncology and similar decision-support tools in oncology;
  \item Watson for Genomics and clinical trial matching systems;
  \item population health / analytics platforms based on acquired datasets;
  \item imaging and radiology decision support initiatives.
\end{itemize}

However, performance in real-world deployments lagged far behind marketing:

\begin{itemize}
  \item The MD Anderson Oncology Expert Advisor (OEA) project, using Watson
technology, reportedly consumed on the order of \$62M and was halted
        after internal review and integration difficulties, including challenges
        interfacing with EMR systems.
  \item Watson for Oncology was criticised for recommending unsafe or
        incorrect treatments in internal tests and for failing to match
        local practice patterns in partner hospitals.
  \item Clinicians often found Watson’s recommendations either obvious or
        impractical, and user engagement remained low.
  \item By 2018–2019, IBM laid off significant staff in Watson Health amid
        market disillusionment.
\end{itemize}

In January 2022, IBM agreed to sell Watson Health’s data and analytics assets
to private equity firm Francisco Partners; these assets were later relaunched
under the brand \emph{Merative}. Watson Health, as originally envisioned, then
effectively ceased to exist.

FRFP reads this arc not merely as a business failure, but as a large-scale probe
of what happens when explicit AI systems are embedded inside complex healthcare
institutions with stringent tacit norms.

\subsection{Explicit and Tacit Spaces at the Platform Level}

\subsubsection*{Explicit artefacts at the platform level}

Across Watson Health products, the explicit artefacts included:

\begin{itemize}
  \item large corpora of medical literature, guidelines, and clinical trial
        registries;
  \item structured claims and administrative data from acquisitions
        (Truven, Explorys, etc.);
  \item prototype knowledge bases and rules encoded with clinical partners;
  \item ML/NLP models and pipelines for ingestion, entity extraction,
        matching, and recommendation;
  \item dashboards and UIs for clinicians and administrators.
\end{itemize}

These artefacts were transformed by Watson’s ingestion, extraction, matching,
and ranking components to produce:

\begin{itemize}
  \item ranked treatment options and genomic interpretations;
  \item patient-trial matches;
  \item population risk scores and cost-prediction dashboards;
  \item imaging decision support outputs.
\end{itemize}

\subsubsection*{Tacit judgement and responsibility}

The tacit side spanned multiple actors:

\begin{itemize}
  \item frontline clinicians’ expertise and pattern recognition;
  \item hospital administrators’ understanding of workflows, incentives,
        and regulatory constraints;
  \item IBM engineers’ and product leaders’ intuitions about AI capability
        and product-market fit;
  \item payers’ and regulators’ tacit standards for evidence and safety;
  \item patients’ expectations and trust.
\end{itemize}

FRFP emphasises that no amount of explicit data and models substitutes for the
tacit clinical and institutional judgment that ultimately governs safe adoption.


\subsection{Boundary Design: From Demos to Clinical Reality}

\subsubsection*{Demos and pilot narratives}

Watson Health’s early narrative leaned heavily on demos:

\begin{itemize}
  \item stylised cases where Watson quickly identified relevant literature or
        treatment options;
  \item inspirational stories of AI helping oncologists keep up with
        exploding literature;
  \item marketing claims about ``fighting cancer'' and transforming care
        through cognitive computing.
\end{itemize}

In FRFP terms, these demos are curated explicit artefacts: carefully chosen
inputs and outputs with limited exposure to the messy realities of clinical
work, local constraints, and safety-critical edge cases.


\subsubsection*{Real-world boundary behaviour}

In deployment, the boundary between Watson outputs and clinical decisions
varied:

\begin{itemize}
  \item Some institutions treated Watson outputs as optional second opinions
        that clinicians could consult or ignore; in these settings, usage
        dwindled when clinicians found little added value.
  \item In others, there was implicit pressure to incorporate Watson into
        workflows due to the cost of the partnership and executive
        expectations, even though clinicians remained formally responsible.
  \item Many deployments never reached stable routine use because system
        integration (EMR, data quality) and clinical trust barriers proved
        too high.
\end{itemize}

From an FRFP perspective:

\begin{itemize}
  \item the intended boundary (AI as support, clinician as decider) was
        not robustly operationalised across sites;
  \item in some contexts, Watson outputs risked boundary capture
        (becoming quasi-authoritative recommendations);
  \item in others, clinicians kept the boundary firmly in tacit space
        by simply not using the tool.
\end{itemize}

\subsection{Protocol-Level Issues in FRFP Terms}

\subsubsection*{Explicit-Space Misalignment: Data, Knowledge, and Tasks}

\paragraph{Heterogeneous, messy data.}
Watson Health attempted to operate on top of highly heterogeneous clinical
data:

\begin{itemize}
  \item EMRs with inconsistent formats, missing data, and noisy free text;
  \item claims datasets optimised for billing rather than clinical meaning;
  \item fragmented, siloed systems that resisted clean integration.
\end{itemize}

In FRFP terms, this means the platform’s “explicit record” of healthcare reality was weaker and noisier than the real clinical situation clinicians were responding to:

\begin{itemize}
\item the platform’s explicit artefacts were not a clean, reliable knowledge base; they were
        a patchwork of partial, inconsistent, and error-prone views;
  \item the mapping from raw EMR data to the structured inputs Watson uses
        is itself a fragile, error-prone transformation step;
  \item many tacit clinical nuances (why a lab was ordered, how a symptom is
        interpreted) never enter the platform’s explicit artefacts at all.
\end{itemize}

\paragraph{Overgeneralised task framing.}
Watson’s Jeopardy! success encouraged a narrative in which:

\begin{quote}
  ingesting lots of text + sophisticated NLP
  leads to good performance in any text-heavy domain, including oncology.
\end{quote}

Critics argue that IBM underestimated how different the tasks are:
Jeopardy! is about pattern-matching and recall; oncology is about causal
and clinical reasoning under uncertainty, with strong safety, workflow, and
ethical constraints.

In FRFP terms:

\begin{itemize}
  \item the explicit transformations Watson applied were suitable for
        question-answering and retrieval in constrained domains, but not for
        approximating tacit oncologist reasoning under uncertainty;
  \item the semantics of clinical decision-making were never properly
represented in $E_{\mathrm{WH}}$; in practice, crucial intermediate
        structure (clinical context, contraindication logic, local pathways,
        and safety-critical nuances) often remained tacit or weakly encoded.
\end{itemize}

\subsubsection*{Boundary and Workflow Design Failures}

\paragraph{Workflow disruption.}
Clinicians in partner hospitals reported that:

\begin{itemize}
  \item Watson’s interface and case-entry workflows were cumbersome and
        disconnected from their usual EMR flows;
  \item using Watson often meant duplicate data entry and extra steps,
        making it unattractive in time-pressured environments;
  \item when Watson outputs duplicated obvious guideline recommendations,
        the extra friction was not justified.
\end{itemize}

From an FRFP angle:

\begin{itemize}
\item the platform’s explicit workflows (data entry, ingestion, recommendation screens, reporting)
        were not aligned with existing clinical and operational workflows and judgment;
  \item boundary events (e.g.\ choosing a treatment plan) were surrounded by
        extra explicit steps that added little tacit value;
  \item the result was low utilisation and erosion of the intended
        explicit/tacit coupling.
\end{itemize}

\paragraph{Lack of prospective evaluation.}
As STAT and others reported, IBM did not initially subject Watson for
Oncology or related tools to rigorous, independent clinical trials before
deploying them commercially.
Instead, validation was often limited, retrospective, or internal.

For FRFP, this is primarily a boundary and endorsement problem:

\begin{itemize}
  \item moving from demo/pilot to real clinical use is a major boundary
        transition requiring strong tacit endorsement from the clinical
        community;
  \item without prospective evidence, the tacit side never fully endorses
        the explicit semantics, leading to persistent distrust and poor
        integration.
\end{itemize}

\subsubsection*{Safe Horizons and Programme-Scale Risk}

Watson Health was not one tool but a portfolio, scaled aggressively:

\begin{itemize}
  \item multiple product lines launched in parallel, across oncology,
        genomics, imaging, and population health;
  \item large, multi-year contracts and acquisitions committed IBM to a
        long horizon before clear evidence of clinical value emerged;
  \item layoffs and write-downs began only after years of limited traction
        and critical reporting.
\end{itemize}

In FRFP’s long-horizon view:

\begin{itemize}
\item Watson Health as a whole is best understood as a long-running, multi-branch programme:
        many product lines, deployments, integrations, and revisions unfolding over years;
  \item safe horizons would require:
        \begin{itemize}
          \item staged pilots with clear success criteria;
          \item caps on deployment scale before independent validation;
          \item explicit re-grounding after each phase (revisiting tacit
                clinical judgments about utility and safety).
        \end{itemize}
  \item Instead, the programme expanded rapidly, with safe-horizon controls
        largely implicit, financial, or reputational (e.g.\ pullback after
        negative press and partner exits).
\end{itemize}

The eventual sale of Watson Health assets can be read as a late, coarse boundary
event: a corporate-level tacit judgment that the explicit portfolio was not
delivering the expected value under real institutional constraints.

\subsubsection*{Institutional Governance and Narrative}

Watson Health lived at the intersection of several institutional operators:

\begin{itemize}
  \item IBM’s corporate governance and investor expectations;
  \item partner hospitals’ internal politics and incentives;
  \item media narratives about AI revolutions;
  \item regulators’ relatively hands-off stance for decision-support tools.
\end{itemize}

FRFP treats these as institutional forces that shape what gets funded, what gets deployed, and what is treated as acceptable evidence for expansion or continued use:

\begin{itemize}
\item IBM’s corporate incentives pushed toward framing Watson Health as a flagship proof of AI leadership,
        encouraging aggressive scaling and optimistic messaging;
  \item partner hospitals faced mixed incentives: prestige and executive sponsorship on one side,
        clinician burden and scepticism on the other;
  \item media narratives amplified both early hype and later disillusionment, shaping public and investor expectations.

\end{itemize}

These operators did not implement an FRFP-style governance layer explicitly:

\begin{itemize}
  \item there was no cross-institutional body formally charged with defining
        admissible uses, boundaries, and safe horizons for Watson Health
        deployments;
  \item governance was emergent, reactive, and often driven by financial or
        reputational shocks (audits, critical articles, contract
        cancellations).
\end{itemize}

\subsection{Counterfactual: An FRFP-Aligned Healthcare AI Platform}

What might an FRFP-compliant analogue of Watson Health have looked like?

\subsubsection*{Narrow, Evidence-Gated Product Strategy}

\begin{itemize}
  \item Start with a small number of carefully scoped tasks:
        e.g.\ information retrieval, literature summarisation, or trial
        matching in well-defined cancer subtypes;
\item treat each as a distinct workflow with its own artefacts, boundary design, safe-horizon limits,
        and evidence gates before expansion;
  \item require prospective or quasi-prospective trials to demonstrate value
        and safety before scaling beyond pilot sites.
\end{itemize}

\subsubsection*{Strong Explicit/Tacit Interface with Clinicians}

\begin{itemize}
  \item Design interfaces that foreground primary sources (guidelines,
        trial reports) and make Watson’s role clearly that of a
        \emph{search-and-structure engine}, not a quasi-physician;
  \item make boundary events explicit:
        treatment decisions, trial enrolment, and care-pathway selection
        always require human endorsement, logged with context;
  \item allow clinicians to easily see when and why Watson’s suggestions
        diverge from local practice or guidelines, and to feed those
        divergences back into the system as explicit signals.
\end{itemize}

\subsubsection*{Safe Horizons for Deployment and Scaling}

\begin{itemize}
  \item Limit initial deployments to small, well-monitored cohorts with
        explicit sunset or review points;
  \item define quantitative and qualitative criteria for expansion
        (e.g.\ concordance rates with expert panels, impact on outcomes,
        usability metrics);
  \item periodically re-ground the system’s tacit model of clinical utility
        via workshops and audits, adjusting or decommissioning pipelines that
        fail to demonstrate value.
\end{itemize}

\subsubsection*{Institutional Governance Across Stakeholders}

\begin{itemize}
  \item Establish joint FRFP-style governance boards with representatives
        from:
        \begin{itemize}
          \item clinicians and nurses;
          \item IT and data governance;
          \item ethics and patient advocacy;
          \item vendor (IBM-style) technical leads.
        \end{itemize}
  \item Charge these boards with:
        \begin{itemize}
          \item setting boundaries (what the AI may and may not do);
          \item defining safe horizons (scale, duration, and scope limits);
          \item deciding when to pause, pivot, or shut down specific
                deployments based on evidence.
        \end{itemize}
  \item Make external communication (marketing, press releases) consistent
        with the governance body’s actual assessment of capabilities and
        validation status.
\end{itemize}

\subsection{Lessons for FRFP and Platform-Scale Healthcare AI}

Beyond the detailed oncology case in Appendix~\ref{app:frfp-watson-oncology},
the broader Watson Health trajectory highlights several FRFP-relevant points:

\begin{enumerate}
  \item \textbf{Platform dreams outpaced explicit semantics.}
        Watson Health attempted to be a general healthcare AI platform
        before the explicit representations of clinical data and tasks were
        mature enough to support that ambition.

  \item \textbf{Demos are not deployments.}
        High-profile proofs-of-concept in curated settings are runs in a
        very narrow slice of real-world conditions.
        FRFP requires attention to how workflows behave at scale, under
        messy, adversarial, and heterogeneous conditions.

  \item \textbf{Safe horizons must be engineered, not assumed.}
        Large, multi-year, multi-product programmes need explicit phase
        boundaries and go/no-go gates, rather than relying on market
        disappointment as the only stopping condition.

  \item \textbf{Institutional operators can amplify both hype and crash.}
        Corporate, clinical, and media institutions updated tacit beliefs
        in ways that first amplified unrealistic expectations, then swung
        to disillusionment.
        FRFP aims to provide structures where those updates are more
        evidence-driven and incremental.

  \item \textbf{Formal architectures help distinguish structural from
        contingent failure.}
        Watson Health’s problems are often framed as ``bad execution'' or
        ``overhype''.
        FRFP shows a deeper pattern:
        incomplete explicit semantics, weak boundary and workflow design,
        absent safe horizons, and diffuse governance---a pattern likely to
        recur in future healthcare AI platforms unless addressed at the
        architectural level.
\end{enumerate}

Under this reading, IBM Watson Health is not merely a failed product line,
but an early, large-scale probe of what happens when a powerful explicit
AI stack is embedded in one of the most complex tacit domains humanity has:
healthcare delivery.
Its rise and fall provide empirical support for FRFP’s insistence on
careful boundary design, safe horizons, and institutional non-automation
in future platform-scale AI deployments.
